% Template created by Karol Kozioł (www.karol-koziol.net) for ShareLaTeX

\documentclass[a4paper,spanish,9pt]{extarticle}
\usepackage[utf8]{inputenc}

\usepackage[T1]{fontenc}
\usepackage{verbatim}
\usepackage{graphicx}
\usepackage{xcolor}
\usepackage{pgf,tikz}
\usepackage{mathrsfs}

\usetikzlibrary{shapes, calc, shapes, arrows, math, babel}

\usepackage{amsmath,amssymb,textcomp}
\everymath{\displaystyle}

\usepackage{times}
\renewcommand\familydefault{\sfdefault}
\usepackage{tgheros}
\usepackage[defaultmono,scale=0.85]{droidmono}

\usepackage{multicol}
\setlength{\columnseprule}{0pt}
\setlength{\columnsep}{20.0pt}

\usepackage[utf8]{inputenc}
\usepackage[spanish]{babel}
\usepackage{eurosym}

\usepackage{graphicx}
\graphicspath{{../img/}}
\usepackage{svg}

\usepackage{hyperref}

\usepackage{geometry}
\geometry{
a4paper,
total={210mm,297mm},
left=10mm,right=10mm,top=10mm,bottom=15mm}

\linespread{1.3}

\newcommand{\samedir}{\mathbin{\!/\mkern-5mu/\!}}

% custom title
\makeatletter
\renewcommand*{\maketitle}{%
\noindent
\begin{minipage}{0.6\textwidth}
\begin{tikzpicture}
\node[rectangle,rounded corners=6pt,inner sep=10pt,fill=blue!50!black,text width= 0.95\textwidth] {\color{white}\Huge \@title};
\end{tikzpicture}
\end{minipage}
\hfill
\begin{minipage}{0.35\textwidth}
\begin{tikzpicture}
\node[rectangle,rounded corners=3pt,inner sep=10pt,draw=blue!50!black,text width= 0.95\textwidth] {\begin{tabular}{cc} \multirow{2}{1cm}{\includegraphics[width=0.15\columnwidth]{header_right}}& \@author \\ & \ies \end{tabular}};
\end{tikzpicture}
\end{minipage}
\bigskip\bigskip
}%
\makeatother

% custom section
\usepackage[explicit]{titlesec}
\newcommand*\sectionlabel{}
\titleformat{\section}
  {\gdef\sectionlabel{}
   \normalfont\sffamily\Large\bfseries\scshape}
  {\gdef\sectionlabel{\thesection\ }}{0pt}
  {
\noindent
\begin{tikzpicture}
\node[rectangle,rounded corners=3pt,inner sep=4pt,fill=blue!50!black,text width= 0.95\columnwidth] {\color{white}\sectionlabel#1};
\end{tikzpicture}
  }
\titlespacing*{\section}{0pt}{15pt}{10pt}


% custom footer
\usepackage{fancyhdr}
\makeatletter
\pagestyle{fancy}
\fancyhead{}
\fancyfoot[C]{\footnotesize \@author \ - \ies}
\renewcommand{\headrulewidth}{0pt}
\renewcommand{\footrulewidth}{0pt}
\makeatother
\usepackage{multirow} % para las tablas


\title{Geometría Analítica}
\author{Departamento de Matemáticas}
\date{2014}
\newcommand{\ies}{IES Pedro Cerrada}



\begin{document}

\maketitle



\begin{multicols*}{2}


\section{Vectores Libres}

Dados dos puntos en el plano (A y B, podemos trazar  una flecha que vaya del primero al segundo. A esta flecha la llamaremos vector (fijo) y se denota $\overrightarrow{AB}$ o $\overrightarrow{v}$.


\begin{tikzpicture}[scale=0.75]
\coordinate (A) at (1,2);
\coordinate (B) at (6,4);

\draw [fill=blue] (A) circle (2pt) node [left] {A};
\draw [fill=blue] (B) circle (2pt) node [left] {B};

\draw [->, red, thick]  (A) -- node[below] {$\overrightarrow{AB}$} (B); 
\end{tikzpicture}



\begin{itemize}
\item \textbf{Módulo}: La longitud del vector
\item \textbf{Dirección}: La recta que contiene al vector y cualquiera de sus paralelas
\item \textbf{Sentido}: El que va del origen al final o su contrario. Viene representado por punta "la cabeza de la flecha"
\end{itemize}
Dos vectores (fijos) son \textbf{equipolentes} cuando tienen el mismo módulo, misma dirección y mismo sentido. Un vector fijo y todos sus equipolentes forman lo que de denomina un \textbf{vector libre}. Un vector libre viene determinado por sus coordenadas:

\section{Coordenadas y módulo de un vector} Un vector se puede ver como el desplazamiento que tenemos que hacer horizontalmente y verticalmente para ir del origen al extremo del mismo. Al desplazamiento horizontal le llamaremos primera coordenada y al vertical, segunda.

\begin{itemize}
\item Dados $A(x_1,y_2),B(x_2,y_2) \to  \overrightarrow{AB}(x_2-x_1,y_2-y_1)$
\item A partir de las coordenadas del punto podremos calcular su módulo. Dados $\overrightarrow{u}(x,y), \to  \left|\overrightarrow{u}\right|=\sqrt{x^2+y^2}$
\end{itemize}




\subsection{Ejemplo}

Determina las coordenadas y el módulo del vector libre cuyo representante es el vector que va de $A(1,1)$ a $B(7,5)$

\begin{tikzpicture}[line cap=round,line join=round,>=triangle 45,x=1cm,y=1cm, scale=0.8]
\draw [color=lightgray,dash pattern=on 1pt off 1pt, xstep=1cm,ystep=1cm] (-0.6129302567150502,-0.43158220601634095) grid (9.010648940148005,6.1);
\draw[->,color=black] (-0.6129302567150502,0) -- (9.010648940148005,0);
\foreach \x in {,1,2,3,4,5,6,7,8,9}
\draw[shift={(\x,0)},color=black] (0pt,2pt) -- (0pt,-2pt) node[below] {\footnotesize $\x$};
\draw[->,color=black] (0,-0.43158220601634095) -- (0,6.1);
\foreach \y in {,1,2,3,4,5,6}
\draw[shift={(0,\y)},color=black] (2pt,0pt) -- (-2pt,0pt) node[left] {\footnotesize $\y$};
\draw[color=black] (0pt,-10pt) node[right] {\footnotesize $0$};
\clip(-0.6129302567150502,-0.43158220601634095) rectangle (9.010648940148005,7.8783927087822985);
\draw [->,line width=1.5pt,color=red] (1,1) -- node[above,fill=white]{$\overrightarrow{v}=\left(6,4 \right) \land \left|\overrightarrow{v}\right|=\sqrt{6^2+4^2}=\sqrt{52}$} (7,5);
\draw [-,line width=2pt,color=blue] (1,1) -- node[below]{$7-1=6$} (7,1);
\draw [-,line width=2pt,color=orange] (7,1) -- node[right]{$5-1=4$}(7,5);

\draw [fill=yellow, color=blue]  (1,1) circle (2pt) node[above] {$A = (1, 1)$};
\draw [fill=yellow, color=blue]  (7,5) circle (2pt) node[above] {$B = (7, 5)$};
\end{tikzpicture}



\section{Operaciones con vectores}
\subsection{Producto de un número por un vector}
\paragraph*{Definición}
Dado $k \in \mathbb{R}$ y $\overrightarrow{u}$ se define $k\cdot\overrightarrow{u}$ como un $\overrightarrow{v}$ que:\begin{itemize}
\item $\left|\overrightarrow{v}\right|=\left|k\right|\cdot\left|\overrightarrow{u}\right|$
\item $ \overrightarrow{v}\samedir\overrightarrow{u}$
\item Mismo sentido que $\overrightarrow{u}$ si $k>0$ o sentido contrario si $k>0$
\end{itemize}Además se cumple que si $\overrightarrow{u}(x_1,y_1)\to k\overrightarrow{u}(k\cdot x_1,k\cdot y_1)$

\subsubsection{Ejemplos}\begin{tikzpicture}[scale=0.75]
\draw [->, red, thick, fill=red]  (1,1) -- node[left] {$\overrightarrow{u}(1,2)$} (2,3); 
\draw [->, blue, thick, fill=blue]  (2,-1) -- node[left] {$2\overrightarrow{u}(2,4)$} (4,3); 
\draw [->, gray, thick, fill=blue]  (6,1) -- node[left] {$\frac{1}{2}\overrightarrow{u}(0.5,1)$} (6.5,2); 
\draw [->, orange, thick, fill=blue]  (8,3) -- node[right] {$-\frac{3}{2}\overrightarrow{u}(1.5,3)$} (6.5,0); 
\end{tikzpicture}

\subsection{Suma y resta de vectores}
\paragraph*{Definición de suma} Dados $\overrightarrow{u}$ y $\overrightarrow{v}$ se define la suma como el vector que   si los ponemos seguidos va del origen del primer vector al extremo del segundo vector. Además se cumple que si $\overrightarrow{u}(x_1,y_1) \ y \overrightarrow{v}(x_2,y_2) \to \overrightarrow{u} +\overrightarrow{v}(x_1+x_2, y_1+y_2)$ 
\medskip 

\begin{tikzpicture}[scale=0.75]
\draw [->, red, thick, fill=red]  (1,0) -- node[left] {$\overrightarrow{u}(1,2)$} (2,2); 
\draw [->, blue, thick, fill=blue]  (2,2) -- node[left] {$\overrightarrow{v}(4,1)$} (6,3); 
\draw [->, orange, thick, fill=blue]  (1,0) -- node[right] {$\overrightarrow{u}+\overrightarrow{v}(5,3)$} (6,3); 
\end{tikzpicture}

\paragraph*{Definición de resta} Dados $\overrightarrow{u}$ y $\overrightarrow{v}$ se define la resta como la suma del primero con el opuesto del segundo. Además se cumple que si $\overrightarrow{u}(x_1,y_1) \ y \overrightarrow{v}(x_2,y_2) \to \overrightarrow{u} -\overrightarrow{v}(x_1-x_2, y_1-y_2)$ 

\section{Punto medio de un segmento}

Dados dos puntos del plano, $A(x_1,y_1)$ y $B(x_2,y_2)$, el punto medio es $M(\frac{x_1+x_2}{2},\frac{y_1+y_2}{2})$. 

\begin{tikzpicture}[scale=0.75]

\tikzmath{\x1 = 1; \y1 =2;
		  \x2 = 6; \y2 =4; 	 
                   %computations are also possible
                    \x3 = (\x1 + \x2)*0.5; \y3 =(\y1 + \y2)*0.5; } 

\coordinate (A) at (\x1,\y1);
\coordinate (M) at (\x3,\y3);
\coordinate (B) at (\x2,\y2);

\draw [fill=blue] (A) circle (2pt) node [left] {$A(x_1,y_1)$};
\draw [fill=blue] (B) circle (2pt) node [right] {$B(x_2,y_2)$};
\draw [fill=blue] (M) circle (2pt) node [above] {$M(\frac{x_1+x_2}{2},\frac{y_1+y_2}{2})$};
\draw [->, red, thick]  (A) -- node[below] {$\overrightarrow{AB}$} (B); 

\end{tikzpicture}

La demostración es sencilla aplicando la propiedad geométrica que cumple el punto medio: $\overrightarrow{AM}=\overrightarrow{MB}$

\subsection{Ejemplo}
\begin{tikzpicture}[scale=0.5]

\tikzmath{\x1 = 1; \y1 = 4;
		  \x2 = 9; \y2 = 0; 	 
                   %computations are also possible
                    \x3 = (\x1 + \x2)*0.5; \y3 =(\y1 + \y2)*0.5; } 

\coordinate (A) at (\x1,\y1);
\coordinate (M) at (\x3,\y3);
\coordinate (B) at (\x2,\y2);

\draw [fill=blue] (A) circle (2pt) node [left] {A(\x1,\y1)};
\draw [fill=blue] (B) circle (2pt) node [right] {B(\x2,\y2)};
\draw [fill=blue] (M) circle (2pt) node [above] {M(\x3,\y3)};
\draw [->, red, thick]  (A) -- node[below] {$\overrightarrow{AB}$} (B); 

\end{tikzpicture}

\section{Puntos alineados}
Dados los puntos $A$, $B$ y $C$ estarán alineados si los vectores $\overrightarrow{AB}$ y $\overrightarrow{BC}$ son colineales, o tienen la misma dirección, y por tanto:
$\exists k\in\mathbb{R} | \overrightarrow{BC}=k\cdot \overrightarrow{AB}$

\begin{tikzpicture}[scale=0.75]

\tikzmath{\x1 = 1; \y1 =1;
		  \x2 = 3; \y2 =2;
		  \k=2; 	 
                   %computations are also possible
                    \x3 = \x2+(\x2 - \x1)*\k; \y3 =\y2+(\y2 - \y1)*\k; } 

\coordinate (A) at (\x1,\y1);
\coordinate (C) at (\x3,\y3);
\coordinate (B) at (\x2,\y2);

\draw [fill=blue] (A) circle (2pt) node [left] {A};
\draw [fill=blue] (B) circle (2pt) node [above] {B};
\draw [fill=blue] (C) circle (2pt) node [right] {C};
\draw [->, red, thick]  (A) -- node[below] {$\overrightarrow{AB}$} (B); 
\draw [->, red, thick]  (B) -- node[below, fill=white] {$\overrightarrow{BC}=k\cdot\overrightarrow{AB}$} (C); 

\end{tikzpicture}

\subsection{Ejemplo}
\begin{tikzpicture}[scale=0.75]

\tikzmath{\x1 = 1; \y1 =6;
		  \x2 = 2; \y2 =5.5;
		  \k=3; 	 
                   %computations are also possible
                    \x3 = \x2+(\x2 - \x1)*\k; \y3 =\y2+(\y2 - \y1)*\k; } 

\coordinate (A) at (\x1,\y1);
\coordinate (C) at (\x3,\y3);
\coordinate (B) at (\x2,\y2);

\draw [fill=blue] (A) circle (2pt) node [left] {$A(\x1,\y1)$};
\draw [fill=blue] (B) circle (2pt) node [right] {$B(\x2,\y2)$};
\draw [fill=blue] (C) circle (2pt) node [below] {$C(\x3,\y3)$};
\draw [->, red, thick]  (A) -- node[below] {$\overrightarrow{AB}$} (B); 
\draw [->, red, thick]  (B) -- node[right, fill=white] {$\overrightarrow{BC}=\k\cdot\overrightarrow{AB}$} (C); 

\end{tikzpicture}

\section{Ecuaciones de la recta}

Podemos definir la recta como el lugar geométrico formado por el conjunto de puntos del plano que a partir de un punto fijo siguen una misma dirección. Dado un punto $P(x_0,y_0)$ y un vector $\overrightarrow{d}(v_1,v_2)$, en la recta $r$ se cumple:

\begin{tikzpicture}[domain=-1:6,>=triangle 45]

\tikzmath{\v1 = 3; \v2 =-1;
		  \x0 = 1; \y0=3;
		  \k=1.5; 
		  \x1 = \x0 +\k*\v1; \y1 = \y0 +\k*\v2;	 
          } 
%\draw[very thin,color=gray] (-0.1,-1.1) grid (3.9,3.9);
\draw[->] (0,0) -- (1,0) node[right] {$x$};
\draw[->] (0,0) -- (0,1) node[above] {$y$};
\draw [color=green,thick] plot(\x,\v2/\v1*\x + 10/3) node [below]{$r$}  ;
\draw [fill=blue] (\x0,\y0) circle (2pt) node [left] {$P(x_0,y_0)$};
\draw [fill=blue] (\x1,\y1) circle (2pt) node [above] {$X(x,y)$};
\draw [fill=black] (0,0) circle (1pt) node [left] {$O(0,0)$};
\draw[->, color=red] (0,0) -- node[right] {$\overrightarrow{OP}$}(\x0,\y0);
\draw[->, color=orange] (0,0) -- node[above] {$\overrightarrow{OX}$}(\x1,\y1);

\draw[->, color=blue] (\x0,\y0) -- node[above] {$\overrightarrow{d}(v_1,v_2)$}(\x0+\v1,\y0+\v2);


\end{tikzpicture}

$$\overrightarrow{OX}=\overrightarrow{OP}+\overrightarrow{PX} $$ 

Como $\overrightarrow{PX}$ y $\overrightarrow{d}$ son colineales:

$$\overrightarrow{OX}=\overrightarrow{OP}+\lambda\cdot\overrightarrow{d} $$ 

\subsection{Ecuación vectorial}
Se obtiene a partir de las coordenadas de la expresión anterior :
$$(x,y)=(x_0,y_0)+\lambda\cdot(v_1,v_2)$$

\subsection{Ecuaciones paramétricas}
Se obtienen separando cada coordenada del expresión anterior:
$$\left\lbrace \begin{gathered}
  x = x_0 + \lambda\cdot v_1 \hfill \\
  y = y_0 + \lambda\cdot v_2 \hfill \\ 
\end{gathered}  \right.
$$
 
\subsection{Ecuación continua}
Se obtienen de la anterior despejando $\lambda$ en cada ecuación e igualando la expresiones:
$$ \dfrac{x-x_0}{v_1}=\dfrac{y-y_0}{v_2}$$

\subsection{Ecuación implícita o general}
Operando y reduciendo la expresión anterior llegaremos a una de la forma:
$$Ax+By+C=0$$

\subsection{Ecuación explícita}
Despejando la $y$ en la ecuación anterior obtendremos
$$y=mx+n$$
donde $m$ es la pendiente y $n$ la ordenada en el origen
\paragraph*{Vector director y pendiente de una recta:} 
Dada una recta $r$ de pendiente $m$ entonces el vector $\overrightarrow{v}(1,m)$ es un vector director de la recta. Y al revés, si $\overrightarrow{d}(v_1,v_2)$ es un vector director de la recta, entonces $m=\dfrac{v_2}{v_1}$ es la pendiente de la recta 

\subsection{Ejemplo}
Dada la recta que pasa por $P(1,3)$ y de dirección la marcada por edl vector $\overrightarrow{d}(3,-1)$ determina la ecuación de la misma en sus diferentes variantes:
\begin{itemize}
\item Ecuación vectorial:
$(x,y)=(1,3)+\lambda\cdot(3,-1)$
\item Ecuaciones paramétricas:
$\left\lbrace \begin{gathered}
  x = 1 + 3\lambda \hfill \\
  y = 3 - \lambda \hfill \\ 
\end{gathered}  \right.
$
\item Ecuación continua:
$\dfrac{x-1}{3}=3-y$
\item Ecuación general:
$3y+x-10=0$
\item Ecuación explícita:
$y=-\dfrac{1}{3}x+\dfrac{10}{3}$
\end{itemize}

\section{Condiciones de paralelismo y perpendicularidad}

\begin{itemize}
\item Para que dos rectas sean paralelas basta con que tengan la misma dirección
\item Dada una recta con vector director $\overrightarrow{d}(v_1,v_2)$, un vector director de las rectas perpendiculares será $\overrightarrow{e}(-v_2,v_1)$ 
\end{itemize}

\section{Ecuación de la circunferencia}

Dada una circunferencia de centro  $P(x_0,y_0)$ y de radio $r$:

\begin{tikzpicture}[domain=-1:6,>=triangle 45]

\tikzmath{\x0 = 3; \y0=2.5;
		  \r=2; 
		  \g=-15;
		  } 



%\draw[very thin,color=gray] (-0.1,-1.1) grid (6,6);
\draw[->] (0,0) -- (1,0) node[right] {$x$};
\draw[->] (0,0) -- (0,1) node[above] {$y$};
%\draw [color=green,thick] plot(\x,\v2/\v1*\x + 10/3) node [below]{$r$}  ;
\draw[fill=blue] (\x0,\y0) circle (2pt) node [left] {$P(x_0,y_0)$};
\draw[fill=blue] (\x0,\y0)+(\g:\r) circle (2pt) node [right] {$X(x,y)$};
\draw[] (\x0,\y0) circle (\r) node [left] {$P(x_0,y_0)$};
\draw [fill=black] (0,0) circle (1pt) node [left] {$O(0,0)$};
\draw[->, color=red] (0,0) -- node[left] {$\overrightarrow{OP}$}(\x0,\y0);
\draw[->, color=orange] (0,0) -- node[right] {$\overrightarrow{OX}$} ($(\g:\r)+(\x0,\y0)$);

\draw[->, color=green] (\x0,\y0) -- node[above] {$\overrightarrow{PX}=r$} ($(\g:\r)+(\x0,\y0)$);
\end{tikzpicture}

Los puntos $X(x,y)$ de la misma  cumplen:

$$\left|\overrightarrow{PX}\right|=r $$


Como $\overrightarrow{OP}+\overrightarrow{PX}=\overrightarrow{OX}$, luego $\overrightarrow{PX}=\overrightarrow{OX}-\overrightarrow{OP}$- Por tanto:

$$\left|\overrightarrow{OX}-\overrightarrow{OP}\right|=r $$
Pasando a coordenadas:
$$ \left|(x-x_0,y-y_0)\right|=r $$

Y por el teorema de Pitágoras obtenemos la ecuación de la circunferencia:

$$ (x-x_0)^2+(y-y_0)^2=r^2 $$

\subsection{Ejemplo}

Determina la ecuación de la circunferencia con centro $P(3,1)$ y radio 3:

%expand(Circle(Point(3,1),3).equation())

$$(x-3)^2+(y-1)^2=3^2$$

$$x^2-6x+9+y^2-2y+1=9$$

$$x^2+y^2-6x-2y+1=0$$

\end{multicols*}

\end{document}
